\kbvexercise{ಅಭ್ಯಾಸ ಪತ್ರಿಕೆ ೧೧}

\vspace{2.5pt plus 3pt minus 0.5pt}\noindent\hspace*{1.2em}\kbvnum{(1)}\hspace{0.55em}ವಿಭಕ್ತಿ ಪ್ರತ್ಯಯಗಳು ಯಾವುವು? ಅವುಗಳ ಉಪಯುಕ್ತತೆಯೇನು?\par

\vspace{2.5pt plus 3pt minus 0.5pt}\noindent\hspace*{1.2em}\kbvnum{(2)}\hspace{0.55em}ವಿಭಕ್ತಿ ಪ್ರತ್ಯಯಗಳ ವಿವಿಧ ಪ್ರಕಾರಗಳನ್ನು ಕ್ರಮಾನುಸಾರವಾಗಿ ಬರೆದು ತಿಳಿಸಿ.\par

\vspace{2.5pt plus 3pt minus 0.5pt}\noindent\hspace*{1.2em}\kbvnum{(3)}\hspace{0.55em}ಸಂಬೋಧನಾ ವಿಭಕ್ತಿರೂಪಕ್ಕೆ ಒಂದು ಉದಾಹರಣೆಯನ್ನು ಕೊಡಿ.\par

\vspace{2.5pt plus 3pt minus 0.5pt}\noindent\hspace*{1.2em}\kbvnum{(4)}\hspace{0.55em}“ಮರ” ಶಬ್ದದ ಏಕವಚನ ರೂಪ ವಿಭಕ್ತಿ ಪ್ರತ್ಯಯಗಳನ್ನು ಬರೆಯಿರಿ.\par

\vspace{2.5pt plus 3pt minus 0.5pt}\noindent\hspace*{1.2em}\kbvnum{(5)}\hspace{0.55em}“ಹಾವು” ಶಬ್ದದ ಬಹುವಚನ ರೂಪ ವಿಭಕ್ತಿ ಪ್ರತ್ಯಯಗಳನ್ನು ಬರೆಯಿರಿ.\par

\vspace{2.5pt plus 3pt minus 0.5pt}\noindent\hspace*{1.2em}\kbvnum{(6)}\hspace{0.55em}“ಹುಡುಗನು ನಂದಿಬೆಟ್ಟವನ್ನು ಹತ್ತುತ್ತಿದ್ದಾನೆ” ಈ ವಾಕ್ಯವನ್ನು ಬಹುವಚನ ರೂಪದಲ್ಲಿ ಬರೆಯಿರಿ.\par

\vspace{2.5pt plus 3pt minus 0.5pt}\noindent\hspace*{1.2em}\kbvnum{(7)}\hspace{0.55em}“ಆಳುಗಳು ತೋಟದಲ್ಲಿ ಹಣ್ಣಿನ ಮರಗಳನ್ನು ಗರಗಸದಿಂದ ಕೊಯ್ದು ನೆಲಕ್ಕೆ ಉರುಳಿಸಿದರು.” ಈ ವಾಕ್ಯದಲ್ಲಿರುವ ನಾಮಪದಗಳನ್ನು ಗುರ್ತಿಸಿ, ಅವುಗಳ ವಿಭಕ್ತಿ ರೂಪಗಳನ್ನು ತಿಳಿಸಿರಿ.\par

\vspace{2.5pt plus 3pt minus 0.5pt}\noindent\hspace*{1.2em}\kbvnum{(8)}\hspace{0.55em}ಈ ನಾಮಪದಗಳ ವಿಭಕ್ತಿ ರೂಪವನ್ನು ತಿಳಿಸಿ :\par

\noindent\hspace*{1.8em}ಕಾಡಿನಲ್ಲಿ,\hspace{0.9em}ಹುಲಿಯನ್ನು,\hspace{0.9em}ಪಿಸ್ತೂಲಿನಿಂದ,\hspace{0.9em}ಕಬ್ಬಿಣಕ್ಕೆ,\hspace{0.9em}ಚೇಳು ರೂಪಾಯಿಗೆ,\hspace{0.9em}ಹೊಲದ,\hspace{0.9em}ಸ್ವಾಮಿಯೇ!\par\vspace{3pt plus 2pt}

\vspace{2.5pt plus 3pt minus 0.5pt}\noindent\hspace*{1.2em}\kbvnum{(9)}\hspace{0.55em}ಈ ನಾಮಪದಗಳಲ್ಲಿ ಪ್ರಕೃತಿ ಪ್ರತ್ಯಯಗಳನ್ನು ಬಿಡಿಸಿ ಬರೆಯಿರಿ.\par

\noindent\hspace*{1.8em}ಸಿಂಹವು ಜನರನ್ನು,\hspace{0.9em}ಮೇಲಕ್ಕೆ,\hspace{0.9em}ಹುಡುಗರಿಗೆ,\hspace{0.9em}ಕಲ್ಲಿನಿಂದ,\hspace{0.9em}ಚಿನ್ನದಲ್ಲಿ,\hspace{0.9em}ಮುಂದೆ,\hspace{0.9em}ಮಕ್ಕಳೇ,\hspace{0.9em}ಮಕ್ಕಳಿರಾ,\hspace{0.9em}ಭೀಮನ.\par\vspace{3pt plus 2pt}

\vspace{2.5pt plus 3pt minus 0.5pt}\noindent\hspace*{1.2em}\kbvnum{(10)}\hspace{0.55em}“ಅವನು” ಶಬ್ದದ ವಿಭಕ್ತಿ ಪ್ರತ್ಯಯಗಳನ್ನು ಏಕವಚನ ಹಾಗೂ ಬಹುವಚನ ರೂಪಗಳಲ್ಲಿ ಬರೆಯಿರಿ.\par
